\section{Related Work}
\label{s:relwk}
%\label{s:discuss}

\myparagraph{System supports for heterogeneous computers.}
There are many related research efforts~\cite{phothilimthana2018floem, seshadri2014willow, gu2016biscuit, baumann2009multikernel, barbalace2015popcorn, gamsa1999tornado, silberstein2013gpufs, kim2014gpunet, silberstein2017omnix, tork2020lynx, brokhman2019gaia, shan2018legoos, Asmussen:2016:MHC:2872362.2872371, nider2021last, pemberton2021restless, 10.1145/3458336.3465273, liu2019offloading, cho2020flick}.
%Distributed OSes, e.g., Multikernel~\cite{baumann2009multikernel}, Omnix~\cite{silberstein2017omnix}, and Popcorn~\cite{barbalace2015popcorn}, are proposed to run a single OS (with multiple kernels) to manage heterogeneous hardware resources; however, it is non-trivial to implement such an OS with all required primitives.
Distributed OSes~\cite{baumann2009multikernel, silberstein2017omnix, barbalace2015popcorn} are proposed to run a single OS (with multi-kernels) to manage heterogeneous devices.
% ; however, it is non-trivial to implement such an OS to support unmodified cloud-native applications. 
Some systems use smart devices to offload specific tasks~\cite{phothilimthana2018floem, seshadri2014willow, gu2016biscuit, kim2014gpunet, silberstein2013gpufs}.
% Flick~\cite{cho2020flick} utilizes advanced hardware features to support applications running on a heterogeneous ISA computer.
Floem~\cite{phothilimthana2018floem} proposes programming abstractions for NIC-accelerated applications.
LeapIO~\cite{li2020leapio} leverages co-processors to offload cloud storage services.
E3~\cite{234944} is a microservice execution platform for SmartNIC-accelerated servers, which can significantly improve the energy-efficiency.
% Molecule~\cite{10.1145/3503222.3507732} proposes a heterogeneous serverless framework by proposing a shim layer between multi-OSes and serverless runtime.
% Compared with these works, 
\os is the first work utilizing network abstraction to achieve dynamic split.
Our experience on \sys has proven the feasibility and efficiency, which is a new breakthrough.

% \sys targets on cloud-native applications, and proposes \os, which implements efficient communication but is still compatible with existing abstractions.
% Moreover, these frameworks did not achieve the \emph{partial offloading} which are necessary to effectively utilize DPU.
%as \os does.


\myparagraph{Optimizing cross-PU communication.}
% SAND~\cite{akkus2018sand} proposes local bus for efficient communication between two functions in the same machine for serverless computing.
Prior efforts~\cite{10.1145/2934872.2934897, 10.1145/2872362.2872401, 10.1145/3230543.3230560, 234944} have explored PCIe/DMA for efficient smartNIC-host communication.
% They usually leverage multiple DMA engines to achieve scatter/gather communication.
GPUnet~\cite{kim2014gpunet} proposes a socket abstraction for GPU programs.
Solros~\cite{min2018s} and SPIN~\cite{bergman2017spin} utilizes P2P DMA to optimize communication.
Pond~\cite{10.1145/3575693.3578835} uses CXL for memory pooling.
CXL-Shm~\cite{10.1145/3600006.3613135} utilizes CXL for efficient DSM and resolves partial failure with reference counting.
\os is the first to provide a network abstraction.
 % for unified abstraction and efficiency. 
% In the contrast, \gsocket targets on the cross-PU networking, and eliminate the TCP/IP stack during the data-path socket operations.
% These libraries can also be enhanced with \gsocket so applications can use \gsocket for local communication, while the user-space network stack for real networking.
% We aim to extend \gsocket with systems like GPUnet in the future to support cross-PU communication between accelerators.

% \TODO{Should also discuss SMC-R and other kernel-bypass network stack design.}
\myparagraph{Related network solutions.}
% There are some related network works.
SocksDirect~\cite{10.1145/3341302.3342071} introduces a user-space socket system,
utilizing a user-mode monitor to manage resources and events.
However, it does not target XPU computers, and still relies on reserved memory buffer.
SMC-R~\cite{rfc7609} supports shared memory-based network over RDMA in kernel.
It can effectively manage resources and support commercial apps.
However, the kernel-centric design incurs costs about switching and copying between user and kernel modes.
IX~\cite{belay2014ix} proposes a dataplane OS that provides both high I/O performance and protection.
Other kernel-bypass networking libraries that also preserve the POSIX API, e.g., Arrakis~\cite{peter2015arrakis}, mTCP~\cite{10.5555/2616448.2616493}, F-stack~\cite{f-stack},
usually utilize efficient user-space TCP/IP implementations and DPDK.
MegaPipe~\cite{10.5555/2387880.2387894} and StackMap~\cite{196184} optimize network stack with zero-copy APIs.
\os follows the line and proposes speculative allocation and UNAPI to balance performance and resource efficiency in a heterogeneous computer.


%
%
%
%
%
%
%
%
%
%%
%
%
%
%%
%
%
%
%%
%
%
%
%%
%
%
%
%%
%
%
%
%%
%
%
%
%%
%
%
%
%%
%
%
%
%%
%
%
%
%%
%
%
%
% The following is old content

\iffalse
\myparagraph{Single kernel.}
% \DD{Single kernel: shared states}
% \DD{Singke kernel + Offloading.}
We name most commercial operating system kernels as \emph{Single kernel}.
It means, the system software relies on a single kernel to manage all the hardware resources and provide a consistent view to users.
Cases include Linux, Windows, macOS, L4, seL4, and many others.
They use \emph{shared memory} to manage internal states (or shared states), e.g., a shared and global list to manage free physical pages.
Some optimizations adopt per-core structures~\cite{veal2007performance, gough2007kernel, boyd2010analysis} to optimize the scalability by avoiding synchronization contention, but do not change their shared memory nature.
Shared memory-based management is suitable for homogeneous architectures, e.g., the case that only CPU is SMP and is the only device with computation abilities, as it is easy to design and implement.
Single kernel systems can use smart devices as normal devices --- treating smart devices as secondary citizens, which wastes the devices.
%However, it is not scalable~\cite{?}\DD{references of scalability issues}.

To better utilize the programmability on smart devices, some single kernel systems~\cite{phothilimthana2018floem, seshadri2014willow, gu2016biscuit} adopt the idea of \emph{offloading}, i.e., they manually offload tasks into the devices.
For example, Seshadri et al. proposed Willow~\cite{seshadri2014willow} in 2014, which is an offloading framework exposing the programmability of smartSSD to applications.
Using Willow, applications can offload tasks to the smartSSD, and communicate with the in-SSD tasks through a remote procedure call mechanism.
Even new programming models\cite{phothilimthana2018floem} and runtimes~\cite{seshadri2014willow, gu2016biscuit, kim2014gpunet, silberstein2013gpufs} have been proposed to enable usage of smart devices, %improve the programmability
offloading requires significant re-design efforts.
As an example, Willow does not support dynamic memory allocation, networking, and other services for in-SSD tasks.
%\DD{more on pros and cons for single kernel + offloading.}
It is almost impossible to achieve general-purpose programming using single kernel systems.

\myparagraph{Multikernel.}
Multikernel is an OS structure that uses multiple kernels to manage hardware resources.
Instead of using shared memory to manage internal states, the states are replicated (i.e., replication) among these kernels and are synchronized explicitly with message passing.
This idea is first proposed for scalability~\cite{baumann2009multikernel}.
%\DD{Trvial debate: shared memory vs. message passing}
%\DD{Improper design choice in multikernel: state replication}
In hindsight, replacing shared memory with message passing to improve scalability is not always resonable~\cite{barbalace2015popcorn, barbalace2017s}.
However, one benefit has been uniformity: multikernel abstraction is suitable for heterogeneous hardware like smart devices-armed architectures.
Multikernel systems are capable of deploying kernels on CPUs and smart devices, and synchronize the states through message.
Therefore, we can achieve general-purpose programming by abstracting a single consistent view for users running on different process units~\cite{barbalace2015popcorn}.

Although multikernel is a workable solution, it is far from the ideal.
One inveterate drawback is the significant synchronization costs caused by state replication.
For example, Barrelfish~\cite{baumann2009multikernel} needs to maintain the consistency of physical memory states among kernels, which incurs synchronizations for each memory allocation operation.
The synchronization costs can lead to unacceptable performance slowdown for smart devices-armed architecture with slow interconnection.
%One inveterate drawback is the significant synchronization caused by state replication, which leads to unacceptable performance slowdown for smart devices-armed architecture with slow interconnection (i.e., Change~\uppercase\expandafter{\romannumeral2}).

\myparagraph{Partitioned kernel.}
The demand to control the synchronization (over interconnection) costs force us to rethink OS states and resources distribution and management.
Intuitively, instead of replication, an OS for smart device-armed architectures should partition OS states and hardware resources to multiple kernels,
and a kernel only communicates with others when using resources managed by others.
This OS kernel is named partitioned kernel.

It is actually, not new, to partition OS states and hardware resources to reduce synchronization costs.
Tornado~\cite{gamsa1999tornado} and K42~\cite{appavoo2007experience} introduce clustered objects to manage state sharing through partitioning and replications.
Corey~\cite{boyd2008corey} proposes a set of new OS abstractions following the principle that let the applications control sharing.
Parakernel~\cite{enberg2019faster} partitions the hardware resources to applications and utilize application-level parallelism to use fast I/O devices.
However, these systems are based on single kernel, which is hard to exploit in-device computing ability.
SemperOS~\cite{hille2019semperos} is a multikernel system partitioning states and hardware resources.
However, it is unawre of smartness of devices, and its design is made based on several strict assumptions: many core architecture, dedicated PUs for kernels, infrequently context switching, new hardware extensions, etc.
\fi
